% ===========================================================================
% carta_sometimiento_raccefyn.tex
% Carta de Sometimiento al Editor en Jefe — RACCEFYN (Sección: Ciencias Físicas)
% ===========================================================================
\documentclass[11pt,a4paper]{article}
\usepackage[spanish]{babel}
\usepackage[utf8]{inputenc}
\usepackage[T1]{fontenc}
\usepackage{amsmath,amssymb}
\usepackage{mathptmx}
\usepackage[margin=2.5cm]{geometry}
\usepackage[hidelinks]{hyperref}
\usepackage{parskip}

\begin{document}

\begin{flushright}
Bogotá, D.C., Colombia, 23 de octubre de 2026
\end{flushright}

\vspace{0.3cm}

\noindent\textbf{Doctor}\\
\textbf{Editor en Jefe y Comité Editorial — Sección de Ciencias Físicas}\\
\textit{Revista de la Academia Colombiana de Ciencias Exactas, Físicas y Naturales (RACCEFYN)}\\
Academia Colombiana de Ciencias Exactas, Físicas y Naturales\\
Bogotá, D.C., Colombia

\vspace{0.4cm}

\noindent\textbf{Asunto:} Sometimiento de artículo de investigación original a la sección de \textbf{Ciencias Físicas}\\
\textbf{Título del manuscrito:} \textit{«Inversión acústica de onda completa mediante redes neuronales informadas por la ecuación Eikonal sobre señales crudas de radiofrecuencia ultrasónica»}\\
\textbf{Título en inglés:} \textit{«Full-waveform acoustic inversion via Eikonal physics-informed neural networks on raw ultrasonic radiofrequency signals»}

\vspace{0.3cm}

\noindent Respetado Señor Editor y miembros del Comité Editorial de Ciencias Físicas:

Por medio de la presente someto a su consideración el manuscrito original titulado \textit{«Inversión acústica de onda completa mediante redes neuronales informadas por la ecuación Eikonal sobre señales crudas de radiofrecuencia ultrasónica»}, de autoría de Yerson Hernando León-Montañez (Facultad de Ingeniería, Maestría en Inteligencia Artificial, Universidad de La Salle, Bogotá, Colombia), solicitando respetuosamente que sea evaluado dentro de la sección de \textbf{Ciencias Físicas} de la \textit{Revista de la Academia Colombiana de Ciencias Exactas, Físicas y Naturales}.

El objeto central de este trabajo pertenece estrictamente al campo de la \textbf{física acústica computacional, la propagación de ondas mecánicas en medios viscoelásticos dispersivos y la solución numérica de problemas inversos no lineales gobernados por ecuaciones en derivadas parciales (EDPs)}. A diferencia de los estudios de procesamiento de imágenes médicas o visión por computadora que operan sobre fotografías en escala de grises ya reconstruidas, la presente investigación aborda el problema inverso físico desde las señales crudas de presión acústica en radiofrecuencia (RF, $1.5\text{--}40\text{ MHz}$). Específicamente, se formula una aproximación geométrica de alta frecuencia basada en la ecuación diferencial parcial Eikonal factorizada ($\lvert\nabla\tau(\mathbf{x})\rvert^2 = c(\mathbf{x})^{-2}$) acoplada con la ecuación de transporte disipativo de Beer-Lambert, las relaciones causales de dispersión de Kramers-Kronig y operadores analíticos diferenciables de formación de haces en el dominio espacio-temporal (\textit{Delay-and-Sum} sobre 75 ondas planas multi-ángulo) y en el dominio del número de onda espacial lateral ($k_x$ mediante transformada de Fourier 2D y corrección por pantalla de fase). Asimismo, se introduce una proyección de ortogonalización de Gram-Schmidt en $L^2(\Omega)$ y un operador diferencial Gaussiano en escala de la primera zona de Fresnel junto con regularización variacional en espacios de Sobolev $H^1$ para suprimir el acoplamiento cruzado (\textit{cross-talk}) inducido por la interferencia coherente de dispersores (\textit{speckle}).

Para demostrar la validez física del formalismo propuesto en dos regímenes experimentales de propagación de ondas de retrodispersión unilateral (\textit{pulse-echo}), se evalúan primero adquisiciones de ondas planas pre-formadas sobre fantomas físicos calibrados de referencia internacional (benchmark CUBDL, logrando un residuo de la EDP Eikonal de $2.8 \times 10^{-5}$, $\text{CNR}_c = 9.30 \pm 0.61\text{ dBr}$ y una ganancia de $+20.3\%$ en el factor de coherencia espacial de Mallart-Fink), y posteriormente campos ultrasónicos focalizados a $40\text{ MHz}$ en medios biológicos heterogéneos (repositorio OASBUD), demostrando la recuperación simultánea y cien por ciento ciega de los mapas continuos de velocidad de fase $c(x,z)$ en $\text{m/s}$ y atenuación espectral $\alpha_0(x,z)$ en $\text{dB/MHz/cm}$.

Declaro formalmente que este manuscrito es completamente original, no ha sido publicado previamente ni se encuentra sometido de manera simultánea a ninguna otra revista nacional o internacional. El documento cumple rigurosamente todas las directrices editoriales vigentes de RACCEFYN: títulos en español e inglés inferiores a 30 palabras y 165 caracteres, título corto de 49 caracteres, resúmenes bilingües inferiores a 250 palabras con seis palabras clave separadas por punto y coma, cuatro figuras y dos tablas principales (total de seis elementos gráficos, inferior al límite de siete permitido por la revista), y 33 referencias verificadas en formato APA con nombres completos de revistas sin abreviar. No existe ningún conflicto de intereses respecto a esta publicación.

Agradeciendo de antemano su amable atención y el tiempo dedicado a la evaluación de este trabajo en la sección de Ciencias Físicas, me suscribo con toda consideración.

\vspace{0.8cm}

\noindent Atentamente,

\vspace{0.9cm}

\noindent\textbf{Yerson Hernando León-Montañez}\\
ORCID: \href{https://orcid.org/0009-0000-4439-8661}{0009-0000-4439-8661}\\
Facultad de Ingeniería, Maestría en Inteligencia Artificial\\
Universidad de La Salle, Bogotá, D.C., Colombia\\
Correos electrónicos: \href{mailto:yeleon36@unisalle.edu.co}{yeleon36@unisalle.edu.co} $|$ \href{mailto:yelemo777@gmail.com}{yelemo777@gmail.com}

\end{document}
